\documentclass[leqno, a4paper, 12pt]{jsarticle}
\usepackage{amssymb}
\usepackage{amsthm}
\usepackage{amsmath}
\usepackage{comment}
\usepackage{url}

\usepackage{enumitem}
%\usepackage{showkeys}


%定理環境 thmはあまり使わずpropをなるべく使う。
%主張は基本的に証明の中で使い、そのため番号を付けない。
%注意環境を付けてみた。
\newtheorem{thm}{定理}
\newtheorem{prop}[thm]{命題}
\newtheorem{cor}[thm]{系}
\newtheorem{lem}[thm]{補題}
\newtheorem{df}[thm]{定義}
\newtheorem{exam}[thm]{例}
\newtheorem{fact}[thm]{事実}
\newtheorem*{claim}{主張}
\newtheorem*{cau}{注意}
\newtheorem*{rmk}{注意}
\renewcommand{\proofname}{証明}
%矢印たち。
%\toは\rightarrowと同じ。\getsは\leftarrowと同じ。同値は\iff
\newcommand{\To}{\Longrightarrow}
\newcommand{\Gets}{\LongLeftarrow}
%黒板太字
\newcommand{\nn}{\mathbb{N}}
\newcommand{\zz}{\mathbb{Z}}
\newcommand{\qq}{\mathbb{Q}}
\newcommand{\rr}{\mathbb{R}}
\newcommand{\cc}{\mathbb{C}}
%無限大
\newcommand{\mgn}{\infty}
%差集合は\setminusで統一する。等号の否定は\neq
%真部分集合は\subsetneqq　偏微分は\partial
%ディスプレイスタイル。\dfracでディスプレイ分数
\newcommand{\dpst}{\displaystyle}

\title{論文メモ
\large{\\--- 関数が肩身の狭い思いをする空間 ---}
}
\author{ナンブ　キトラ}
\date{\today}
\begin{document}
\maketitle

本棚を整理したいので，
溜まっている論文のメモを書いて未来への手紙とすることにする．

以下の論文はおそらく，位相空間であって，その上の実連続関数が定数になってしまうものの例を調べようとして集めた論文だと思う．
\section{論文たちの簡易な説明}
\textbf{神話．}
むかしむかし，どれくらい昔かというと
ソビエト連邦という今ではもう存在しないものが出来たてほやほやの時くらいに，
ウリゾーンという数学者がいて(まあ夭折しちゃうわけだけれども・・・)，
ジェネトポのこといろいろ研究していて，
正規空間の上には実連続関数がたくさんあるということを発見したわけ．これがウリゾーンの補題というやつ．どれくらい豊富かというと交わらない二つの閉集合を実連続関数を使って分離できるくらいある．
そのような関数が豊富にある空間は最近だと正規よりちょっと弱い
完全正則空間(completely regular)という
ものが典型例だとみなされていると思う．
完全正則空間というのは，
任意の一点とそれを含まない閉集合
を関数で分離できる程度には実連続関数が
豊富ある空間で，
実連続関数の情報だけから元の位相を復元できるし，
一様付け可能性と同値だったりするしで
かなり便利にいいように扱われていると思う．
とにかく空間を与えた時にその上に豊富に連続関数が
あるのかどうか気になるわけなのだ．

ウリゾーンは
論文
\cite{MR1512258}
で連結ハウスドルフ可算空間を構成したようだ．
しかしながらドイツ語が読めないので
これは
\cite{MR0017527}
の受け売りである．
連結ハウスドルフ可算空間$X$を考えて
実連続関数
$f\colon X\to \rr$を考えると
$f(X)$
は$\rr$
の連結な可算空間でなければならないが
そのような空間は一点集合しかないので
$f$は必ず定数になってしまう．
このように連結ハウスドルフ可算空間
というのは関数が定数だけになる仕組みが
大変分かりやすい空間になっている
(こういう話は\cite{Den1}にも書いた)．
ところで
連結正則な可算空間は存在しないことに注意しよう．
なぜなら
可算空間は常にリンデレフでリンデレフ正則
ならば正則空間になってしまうので
それこそウリゾーンの補題から関数が豊富に存在することが
示せてしまう．
というわけでウリゾーンの補題から始まる数学の歴史は
まあたくさんあるんだけど，
とにかくそのうちの一つは関数が豊富にある空間を扱うような話に発展してそれは完全正則空間の話で，
そしてまた一つは全く正反対の話，つまり関数が定数しかない空間は
どのくらい複雑な位相的性質を許容できるのか？という話に発展したということである．今回の論文の束は
そういう関数が肩身の狭い思いをする空間のお話である．
このような空間はコンパクトではない
pseudo-compact 空間にもなっている．
Pseudo-compactについては
\cite{MR3822332}などを参照のこと．
ただしこの本では空間は全て完全正則と仮定されている
(それはそう)のでこの文章とは少々趣が異なる．

\subsection{Hewitt}
先にも登場したが，
Hewittの論文
\cite{MR0017527}
では
Urysohnの次の二つの問題に答えている．

1. その上の実連続関数が常に定数になるような正則空間は存在するか？

2. 連結空間の最小濃度はどうなるか？特に異なる2点について
これら2点の交わらない二つ閉近傍で分離できるような(このような空間はUrysohn spaceと呼ばれている)
連結空間の最小濃度はどうなるのか？

Hewittはこれらの問いに完全に答えた．
正則基数$\kappa$を与えたときに(
ここでの正則は空間の正則性と全然意味が違うのでややこしい)，
濃度が$\kappa$であるような
正則空間であってその上の実連続関数が
定数になるような空間は存在する．
次に
連結Urysohn 空間の最小濃度は可算であることも証明している．
つまり連結Urysohn 可算空間が存在する．
しかしながらHewittの構成はなんというか
なんだか難しいと思う．

\subsection{チェコ語の論文：読めない}
論文
\cite{novak1948regular}
では
正則空間であって
その上の実連続関数が常に定数になるような
ものを構成している．しかしながら
この論文はチェコ語で書かれており，
とりあえず私には読めない．
一番後ろに英語のサマリーが
載っているのでここだけ読めると思う．
ちなみにこのサマリーによると
Hewittの論文の発表の前の1948年に空間の構成を得ていたが
1946年という年の終わりにHewittの論文のゲラ刷り?が出回ったとか，
Hewittの論文は1945年にsubmitされてて，
結果として一年違いで
同じ問いの解答
を発見したことなどが書いているが，
異なる構成だし発表する価値があるだろうといったことが
誠実に書かれている．
優先権というものがありますからこういうことはちゃんと
書いとかないといろいろ怒られますからね．
しかしまあ，このお二方の構成はちょっと難しいかもしれない．


\subsection{ジョーンズのお仕事とカラクリ (Jones Machine)}
Jones 
\cite{MR0093757}
は
連結，局所連結，
完備なMoore空間$\lambda_{\infty}$であって
ある点$p$では完全正則性が破綻するような空間を構成した．
この空間の構成はJones Machine
と呼ばれるものの萌芽だと思うが
ちょっとよくわかんない．

Jones Machine
(ジョーンズの機械/ジョーンズの機構/ジョーンズのカラクリ)
について説明をしよう．
Jones Machine という名前は
M. E. RudinによるJonesの仕事の解説論文
\cite{MR1617573}
に現れている．
元を辿るとこの
ジョーンズの機構は
論文
\cite{MR0413044}
の構成だと思う．
一体ジョーンズのカラクリとはなんなのか説明しよう．
この論文
\cite{MR0413044}
では
何か位相的性質
$\mathcal{P}$
をもつ非正規空間$X$
が存在した場合に
その位相的性質$\mathcal{P}$
を持つような非完全正則空間を構成する方法が与えられている．もちろん性質$\mathcal{P}$はなんでもいいというわけではないが，とにかくこの構成法が
\textbf{Jones Machine(ジョーンズの機械/ジョーンズの機構/ジョーンズのカラクリ)}
と呼ばれている．
ざっくり解説すると非正規空間$X$の開集合で分離できない
二つの閉集合$H$と
$K$を固定する．そして$X$の可算個のコピーを考える．それらを$X_{i}$と書きその部分集合は下つきの添字$i$をつけることにする．
そして$X_{1}$と$X_{2}$は$H_{1}$と
$H_{2}$を同一視して貼り合わせる．
$X_{2}$と
$X_{3}$は$K_{2}$と
$K_{3}$を貼り合わせる．
そして無限遠点
$\infty$
を付け加えると
$H_{i}$や$K_{i}$が関数で分離できないことが伝播して
無限遠点$\infty$
と閉集合$H_{1}$を分離する関数が存在しないことがわかるというカラクリである．それで実はこのように商を使って構成した空間が正則性を保つというのは位相空間論のいい話の一つ(フォークロア)だったと思う．そういう商の話は大体ブルバキに載ってる．
またこの論文の問題意識はS-spaceやL-spaceという話に
繋がったらしいがここら辺はよくわかんない．


また同じ論文\cite{MR0413044}でJonesは自然数ではなく整数で添字づけして
両端にそれぞれ一つずつ合計二つの無限遠点
$\infty_{-}$
と
$\infty_{+}$を付け加えた構成をすると，
その空間の上の実連続関数$f$について，
$H_{i}$と
$K_{i}$が関数で分離できないということが伝播して
$f(\infty_{+})=f(\infty_{-})$
が成り立つことを注意している．
このように関数が特定の2点で定数になるような空間が
あると，えげつない貼り合わせで
正則空間でその上の実連続関数が定数になるものが構成できる．これから出てくるサブセクションにもそういう構成が紹介されると思う．





論文
\cite{MR1802673}
はJonesの追悼論文だと思うが，
Jonesの業績に関していろいろ説明してくれている．

\subsection{Moore spaceってなんだっけ}
Armentrout
は
論文
\cite{MR0120615}
で
Moore spaceであって
その上の実連続関数が
定数になるものを構成した．
Moore spaceというのは
正則空間であってなんか可算個の被覆の族があって
いい感じの条件を満たすものである．
とにかく正則空間である．

しかしながらこの空間は
Jonesの空間
\cite{MR0093757}
をさらにいじくって構成しているのでなかなか難しい．


\subsection{一般の連続関数が定数になる空間}
Herrlichの論文
\cite{MR0185565}
はドイツ語で書かれているので
英訳されたものをサラッと読んだのだが，
なんかすごい定理が証明されている．
\begin{thm}[\cite{MR0185565}]
$Y$を位相空間とする．この時以下は同値である．
\begin{enumerate}[label=\textup{(\alph*)}]
\item 
$Y$
は$T_{1}$である．
\item 
正則空間
$X$が存在して連続関数
$f\colon X\to Y$
は常に定数になる．
\end{enumerate}
\end{thm}
証明はなんかアレフナンバーを使うので難しく見える(かも)．


\subsection{クラスルームノート}
論文\cite{MR1535290}
では正則だが完全正則ではない空間の例が構成されている．
それは平面上の部分集合に変な位相を入れて仮想的な
二つの無限遠点$p_{+}$と
$p_{-}$からなる空間でこの空間の上の
実連続写像は常に
$f(p_{+})=f(p_{-})$
を満たすようなのでここから
完全正則性を満たさないことがわかる．
こういう空間のコピーをたくさん用意して任意の2点に対して
その2点と
$p_{+}$と$p_{-}$を同一視してくっつけるという操作を可算回繰り返すと原理的にはその上の実連続関数が定数になるものが構成できる．
はてなブログ
\cite{Concon1}にも説明が
あると思う．
こういう話は
\cite{MR0271902}
に載っている．
特に任意の正則空間$Z$を
正則だがその上の実連続関数が
定数になるような空間
$Q(Z)$に位相的に埋め込めることも示している．
論文
\cite{MR0312462}
もそういう話．

\subsection{Younglove}
ヤングラヴ
は
\cite{MR0248741}
で局所連結な
完備Moore空間
でその上の実連続関数が
定数に限るものを構成した．
ちょっと構成はよくわかんなかった．


\subsection{80年代以降}
ここからは80年代以降の話になる．

Mysior は
論文
\cite{MR0601748}
で正則だが完全正則でない
空間を構成した．
その台集合
は上半平面に一点余計な点$a$を付け加えた空間である．
変な位相を入れるとそういう空間になるのである．
さらにその空間に
点$b$を付け加えていい感じに近傍を定義すると
その拡張された空間のについては実連続関数は必ず
$f(a)=f(b)$
を満たしてしまうのである．
そして上で述べた話などから
正則でその上の実連続関数が定数になる空間を作れるので
結局Urysohnの問いにある空間の簡単な構成が得られたことになる．
今現在において正則関数であってその上の実連続関数が定数になってしまうものを構成したい場合には
Mysior の
論文
\cite{MR0601748}
の空間をフォローして，
そういう空間を使う構成の
\cite{MR0271902}
や
\cite{MR0312462}
を適用すると一番手っ取り早いと思う．



論文
\cite{MR0614896}
はなんかよくわかんなかった．
ほんとごめんなさい．
正則性とコンパクト化についての研究らしい．

論文
\cite{MR0712857}
は
Jonesの論文
\cite{MR0413044}
の発展である．

論文
\cite{MR0749120}
は位相空間$X$とその完全正則化
について直積に関する
可換性は成り立たないという論文である．
位相空間の完全正則化というのは$X$のコゼロ集合から
生成される位相をもつ空間である．
多分マクレーンの``圏論の基礎''に書いてある別に完全正則じゃなくても
ストーンチェックコンパクト化は考えられるよ
みたいな話の埋め込み(埋め込みじゃないけど)から誘導される位相空間だと思う．多分なんとか随伴関手とかになっていると思う．おそらくそういう話だと思う．ちょっとそこまで詳しくわかってない．

論文
\cite{MR0732053}
はタイトルの通り
ハウスドルフ空間
$Y$を与えると
非自明なMoore空間$X$であって
連続関数$f\colon X\to Y$
はつねに定数になるものを構成している．
Herrlichの論文
\cite{MR0185565}
の発展と見做せる．
Herrlichの論文
\cite{MR0185565}
がよくわからなくていろいろ探してこの論文に行き着いたような記憶がある．

これはもう90年代の論文だが
論文\cite{MR1489200}
でTzannes
はハウスドルフな可算コンパクト
空間であってその上の実連続関数が常に定数になってしまう例を構成している．
これはコンパクトハウスドルフなら正規でその上の連続関数は豊富にあるという話のコンパクト性を可算コンパクトに弱めると豊富さが損なわれるという話だと思う．

これ以外のものについてはこれらの論文の参考文献や
これらの論文を引用している論文をGoogle Scholarなどで
調べれば大体出てくるのではないかと思われる．


\section*{参考文献たち}

\renewcommand{\refname}{はてなブログ}
	
\begin{thebibliography}{99}
\bibitem[電波通信1]{Den1}
はてなブログ電波通信の記事
「連結ハウスドルフ可算空間」, 
https://concious4410.hatenablog.com/entry/2016/12/16/193128

\bibitem[狐1]{Concon1}
はてなブログ CommonNoun's diary の記事
「Urysohnの距離化定理とその周辺 その２」, 
https://commonnoun.hatenadiary.jp/entry/2016/12/30/022732

\end{thebibliography}


\renewcommand{\refname}{参考文献}

%READ!
%myplain is the same to amsplain style. 
\nocite{*}
\bibliographystyle{myplaindoidoi}
\bibliography{const.bib}



\end{document}